\scanpage{164}
\mspage{1}
% archive sheet number: 163

\unsure{VI|VII} \uline{Groupe de monodromie et groupe fondamental} $(*)$.

\deleted{\illegible{1}} Dans les exposés précédents, nous avons étudié, \deleted{l'ac} pour un schéma de type fini $X$ sur $K$ (corps des fractions de l'anneau de val. \unsure{discr.|discr} str. local $V$), l'action du groupe d'inertie $I=\operatorname{Gal}(\bar K/K)$ sur les gpes de cohomologie de $X_{\bar K}$. Il y a lieu également d'étudier l'action de $I$ sur \deleted{les gpes} d'autres invariants de nature topologique \deleted{\unsure{associés}} de $X_{\bar K}$. Dans le présent exposé, nous nous intéresserons surtout, à l'action de $I$ sur le groupe fondamental ${\unsure{\pi|\Pi}}_1(X_{\bar K},\unsure{\xi|\zeta})$ (relativement à un pt géom $\xi$ de $\bar K$). \deleted{Nous désignerons par $\pi_1^{(\unsure{p|p'})}(X_K,\xi)$ le}

1. \deleted{\unsure{Résultats}}\uline{Énoncé des résultats}.

1.1. Pour un groupe profini $\pi$, nous désignerons par la suite par $\pi^{(p')}$ le plus grand quotient pro-fini de $\pi$ premier à $p$ \inserted{(\unsure{i.e.|soit} $p'=\mathbb{P}-\{p\}$, et)} (où $p$ est l'exp. car. de $k=k(V)$). Donc si $p=1$, \inserted{i.e. $k$ de car. nulle,} on a $\pi^{(p')}=\pi$, sinon $\pi^{(p')}$ est le quotient de $\pi$ par le ss-gpe invariant fermé \uline{engendré} par \unsure{un|son} $p$-ss-groupe de Sylow, i.e. le quotient de $\pi$ par le ss-gpe engendré par les $p$-ss-groupes de Sylow de $\pi$.
% margin note written sideways
\margin{$(*)$ Le présent exposé n'est pas nécessaire pour la compréhension du reste du séminaire.}

La formation de $\pi^{(p')}$ est évidemment fonctorielle en $\pi$, et transforme autos intérieurs de $\pi$
